\documentclass{beamer}
\usepackage{hyperref} % URLs
\title[Critical Thinking]{Science, truth, and honesty }
%\subtitle[short]{full}
\author{Andy J. Wills}

\date{}

\begin{document}
%%%%% All images from wikimedia commons.
\frame{\titlepage}


\begin{frame}{Topics}
	\begin{itemize}
		\item What is truth? 
		\item What is science?  
	\end{itemize}
	\centerline{
	\includegraphics[width=.3\textwidth]{pics/truth.jpg}
	\includegraphics[width=.3\textwidth]{pics/science.png}
      }

\end{frame}

\begin{frame}{Truth as a property of claims}
	\begin{itemize}
		\item Statements about truth
		\begin{itemize}
			\item ``Andy is in Plymouth'' $\to$ TRUE 
			\item ``Andy is on the Moon'' $\to$ FALSE 
		\end{itemize}
		\item Statements about the limits of our knowledge
		\begin{itemize}
			\item ``Andy was born in the month of June'' $\to P(true) = 1/12$ 
			\item ``Andy has \pounds 3.57 in his pocket'' $\to$ UNKNOWN
		\end{itemize}  
	\end{itemize}
\end{frame}

\begin{frame}{Subjective and objective claims}
	\begin{itemize}
		\item \textbf{Subjective claims} are those whose truth differs for different people.
		\item ``My favourite colour is blue''
		\item \textbf{Objective claims} are those whose truth value is not affected by who says it.
		\item ``Charlotte Smith's favourite colour is blue''
		\item The term \emph{subjective} is widely misused.
		\item ``I think chocolate tastes better than cabbage'' is subjective.
		\item ``Chocolate tastes better than cabbage'' is objective, and is amenable to scientific test.
	\end{itemize}
\end{frame}

\begin{frame}{Testing a claim}
Online poll 1 ... 
\end{frame}

\begin{frame}{Vague claims are subjective}
\begin{itemize}
\item ``Smoking is wrong'' - Subjective.
\item ``Smoking increases life expectancy'' - Objective.
\item Note: Claims do not need to be correct to be objective.
\end{itemize}
\end{frame}

\begin{frame}{Objective or subjective?}
Online poll 2 ...
\end{frame}

\begin{frame}{Scientific claims}
	\begin{itemize}
		\item Science is the process of making scientific claims...
		\item ...and determining whether those claims are true or false.
		\item Scientific claims are \textbf{objective} rather than \textbf{subjective}.
		\item Scientific claims are those whose truth can, at least in principle, be clearly determined.
		\item Scientific claims are largely \textbf{descriptive} rather than \textbf{prescriptive}.
	\end{itemize}
\end{frame}

\begin{frame}{Descriptive versus prescriptive claims}
	\begin{itemize}
		\item \textbf{Descriptive claims} say something about how the world is, was, or will be:
		\begin{itemize}
			\item ``Plymouth University was once a polytechnic''
			\item ``Reaction times slow as people age''
			\item ``The UK's Gross Domestic Product will grow by 3\% next year''
		\end{itemize}
	\end{itemize}
\end{frame}


\begin{frame}{Descriptive versus prescriptive claims}
	\begin{itemize}
		\item \textbf{Prescriptive claims} say something about how how things \emph{should} be:
		\begin{itemize}
			\item ``Abortion should be made illegal''
			\item ``People should not drink and drive''
		\end{itemize}
		\item How to examine prescriptive claims scientifically:
		\begin{itemize}
			\item Look for a descriptive claim that might support the prescriptive claim:
			\item \emph{Anti-abortion}: ``feotal pain receptors have developed by eight weeks gestation''
			\item \emph{Pro-abortion}: ``foetal pain receptors are not connected to the brain until 20 weeks''
			\item \emph{Anti-drink-driving}:  ``Risk of car accidents doubles at 80mg/100ml blood alcohol (UK drink-driving limit)''
		\end{itemize}
	\end{itemize}
\end{frame}
			

\begin{frame}{Descriptive versus prescriptive claims}
	\begin{itemize}		
		\item Science is not about avoiding societally difficult questions; nor is it exclusively about societal impact.
		\item It's about making claims that can be reliably examined.
		\begin{itemize}
			\item Some of those claims have direct societal impact (drink driving claims)
			\item Other claims change the way we view ourselves and the world in the longer term.
		\end{itemize}
	\end{itemize}
\end{frame}

\begin{frame}{Absolute versus contextual claims}
	\begin{itemize}	
		\item 	\textbf{Absolute claims} are invariant. They hold always. Their truth value is not conditional on circumstances. They are not conditional on time or place.
		\item \textbf{Contextual claims} hold under a defined set of conditions.
		\begin{itemize}
			\item ``The leadership positions that women occupy are less promising than those of their male counterparts''
			\item Not intended to be an absolute claim.
			\item If the claim is true now, but no longer true in 20 years, this does not undermine the truth value of the original claim.
		\end{itemize}
		\item Scientists generally wish to make claims that are as context-independent as possible; otherwise it is hard to make further predictions.
		\begin{itemize}
			\item ``The leadership positions that women in 2003 in UK FTSE100 companies occupy are less promising than those of their male counterparts''
		\end{itemize}
	\end{itemize}
\end{frame}

\begin{frame}{Observable, measurable states}
	\begin{itemize}
		\item Scientific claims are based on observable, measurable states
		\item Objective, descriptive claims are not always measurable.
		\item ``Impulsive people are more likely to be criminals''
		\begin{itemize}
			\item Being a criminal is a state that is observable and measurable.
			\item Impulsivity is a vague concept that must be translated into something measurable.
			\item e.g. Barratt Impulsivity Scale.
		\end{itemize}
		\item ``People with one or more criminal convictions score higher on the BIS than people without a conviction''
	\end{itemize}
\end{frame}

\begin{frame}{Independent replication}
	\begin{itemize}
		\item Scientific claims must be expressed in such a way that they permit independent replication.
		\begin{itemize}
			\item ``Willsian Therapy reduces depression''
			\item Wills is the only person who can perform Willsian therapy
			\item The claim is not scientific because any attempt to assess its truth value would have to involve Wills and hence could not be independently replicated.
		\end{itemize}
	\end{itemize}
\end{frame}

\begin{frame}{Scientific claims}
	\begin{itemize}
		\item \textbf{O}bjective
		\item \textbf{D}escriptive 
		\item \textbf{A}ppropriately context-independent
		\item \textbf{T}rue or false	
		\item \textbf{B}ased on observable measurable states
		\item \textbf{I}ndependent replication
		\item \textbf{F}alsifiable
	\end{itemize}
\end{frame}

\begin{frame}{Activity 1: Scientific claims}
  Come up with, or find, three psychology-relevant claims that are:
  \begin{enumerate}
    \item Largely unscientific (0-1 ODA.BIF)
    \item Somewhat scientific (2-3 ODA.BIF)
    \item Largely scientific (5-6 ODA.BIF)
  \end{enumerate}
  In your group:
  \begin{itemize}
  \item Introduce yourselves
  \item Pick one favourite psychology topic each
  \item Decide which is going to be 1,2,3 scientific
  \item Write, or find online, a claim for each
  \end{itemize}
  One source of claims: \url{https://nobaproject.com/browse-content}  
\end{frame}

\begin{frame}{Activity 1}
  \begin{itemize}
  \item Claims you came up with.
  \item General Q \& A for part 1.
  \end{itemize}
\end{frame}

\begin{frame}{Diderik Stapel}
	\includegraphics[width=.9\textwidth]{pics/stapel.jpg}

	Dutch social psychologist - admitted to inventing data
	
	See also - Dirk Smeesters - Dutch social psychologist - ``cherry picked'' data.
	
\end{frame}

\begin{frame}{Science and dishonesty}
	\begin{itemize}
		\item Dishonesty
		\begin{itemize}
			\item Partially reporting your results (if your intention is obfuscation). 
			\item Choosing a form of data analysis because it gives you the result you want.
			\item Publishing the same data more than once.
			\item Making up data!
		\end{itemize}
		\item Dishonesty gets in the way of reliably evaluating claims.
	\end{itemize}
\end{frame}

\begin{frame}{Failures to replicate}
	\centerline{\includegraphics[width=.2\textwidth]{pics/ap.jpg}}

	Ap Dijksterhuis. Dutch social psychologist - his ``intelligence priming'' effects fail to replicate. 
	
	\url{http://www.nature.com/news/disputed-results-a-fresh-blow-for-social-psychology-1.12902}
	
	See also - John Bargh - US social psychologist - his ``age and walking pace'' results fail to replicate
	
\end{frame}

\begin{frame}{Reproducibility project}
	\centerline{\includegraphics[width=.8\textwidth]{pics/nosek.png}}
	\url{http://www.sciencemag.org/content/349/6251/aac4716}
\end{frame}

\begin{frame}{Activity 2: Does it replicate?}
  Ludmer et al. (2011). ``Uncovering camouflage: amygdala activation predicts
  long-term memory of induced perceptual insight''

  In your group: 
  \begin{enumerate}
    \item Use Google Scholar to find Ludmer et al. (2011)
    \item Click ``cited by'' to get list of citations
    \item Split pages of citations between you
    \item To get freely-available versions of papers, you may need to click ``all X versions``
    \end{enumerate}

\end{frame}

\begin{frame}{Activity 2: Does it replicate?} 
  On your pages, search for replications:
  \begin{itemize}
  \item Following words in title, journal name or abstract means you can skip ``review``, ``meta-analysis``, ``lecture``, ``theory article``, ``discussion article`` - these are discussions, not attempts to replicate.
  \item No Results section - not an experiment - skip
  \item No pictures of brains - not a neuroscience study - skip
  \item Not studying humans - skip
  \item Not written in English - you can skip for this exercise
  \item Not a replication - e.g. studies speech, musical creativity
  \item If it might be a replication i.e. does an aha! experiment or something like it - are the results the same? Specifically, is activation of the amygdala demonstrated?
  \end{itemize}
\end{frame}

\begin{frame}{Activity 2}
  \begin{itemize}
  \item Any potential replications? Did they work?
  \item General Q \& A for part 2.
  \end{itemize}
\end{frame}

\begin{frame}{RStudio login}
  \begin{itemize}
    \item Check you have received email with login details
    \item Check your login works: rstudio.plymouth.ac.uk
    \item Any issues, contact Tech office urgently to resolve: tech.psy@plymouth.ac.uk
    \end{itemize}
\end{frame}

\begin{frame}{Further materials}
  \begin{itemize}
  \item The notes accompanying this lecture
  \item NOVA chapter: \url{ https://nobaproject.com/modules/the-replication-crisis-in-psychology}
  \end{itemize}
  \vspace{12pt}    
      \tiny
These slides are licensed under a Creative Commons Attribution-ShareAlike
4.0 International Licence. Last update: \today
    
\end{frame}


\end{document}



%%% Local Variables:
%%% mode: latex
%%% TeX-master: t
%%% End:
